\documentclass{article}
\usepackage{graphicx}
\usepackage{amsmath}
\usepackage{amsfonts}
\usepackage{bookmark}
\usepackage{booktabs}
\usepackage{exercises}
\graphicspath{ {./images/10.1} }
\usepackage{indentfirst}
\usepackage{amsthm}
\usepackage{yhmath}
\usepackage{geometry}
\usepackage{placeins}
\usepackage{csquotes}
\usepackage{pgfplots}
\newcommand*{\oldneg}{\mathord{\sim}}
\newcommand\showdiv[1]{\overline{\smash{)}#1}}
\newcommand\Mydiv[1]{%
\smash{\raise.3ex\hbox{$\big)$}}\mkern-8mu
        \overline{\enspace\strut#1}}
\setlength{\tabcolsep}{5pt}
\theoremstyle{definition}
\newtheorem{definition}{Definition}

\setlength{\parindent}{.5in}
\setlength{\parskip}{12pt}
\renewcommand{\arraystretch}{2}
\geometry{margin=1in}
\newcommand{\?}{\stackrel{?}{=}}
\pgfplotsset{width=7cm,compat=1.18}\usepgfplotslibrary{statistics}

\title{10.2 Answers}
\begin{document}
\maketitle

\section*{1}
The graph of $y=\cos 2(x-3)$ is the graph of $y=cos 2x$ translated two units to the right.

\section*{2}
The graphs $y=\tan x$ and $y=-\tan x$ are reflections of eachother over the $x$-axis. In addition, they have all of the same $x$-intercepts.

\section*{3}
E

\section*{4}
C

\section*{5}
D

\section*{6}
B

\section*{7}
A

\section*{8}
F

\section*{17}
We are given the function $y=4\sin\frac{1}{3}\left(x+\frac{\pi}{2}\right)$.
We know that the horizontal shift is $h=-\frac{\pi}{2}$ and the vertical shift is $k=0$. Also, the period is $6\pi$ and the amplitude is 4. Thus the $x$-intercepts are
\[ x=\frac{6\pi}{2}-\frac{\pi}{2}+6\pi n=\frac{5\pi}{2}+6\pi n\vee x=6\pi-\frac{\pi}{2}+6\pi n=\frac{11\pi}{2}+6\pi n \] where $n$ is an integer and the maximum values are
\[ (\frac{1}{4}\times 6\pi-\frac{\pi}{2}+6\pi n,4)=(\pi+6\pi n,4) \] where $n$ is an integer and the minimum values are
\[ (\frac{3}{4}\times 6\pi-\frac{\pi}{2}+6\pi n,-4)=(4\pi+6\pi n,-4) \] where $n$ is an integer:\newline

\section*{18}
We are given the function $f(x)=\cos\left(x-\frac{\pi}{8}\right)+2$.
We know that the horizontal shift is $h=\frac{\pi}{8}$ and the vertical shift is $k=2$. Also, the period is $2\pi$ and the amplitude is 1. Thus the midline is $y=2$, and the points on it are
\[ x=\frac{1}{4}\times 2\pi+\frac{\pi}{8}+2\pi n=\frac{5}{8}\pi+2\pi n\vee \]
\[ x=\frac{3}{4}\times 2\pi+\frac{\pi}{8}+2\pi n=\frac{7}{8}\pi+2\pi n \] where $n$ is an integer and the maximum values are
\[ (2\pi+\frac{\pi}{8}+2\pi n,1+2)=(\frac{17\pi}{8}+2\pi n,3) \] where $n$ is an integer and the minimum values are at
\[ (\pi+\frac{\pi}{8}+2\pi n,-1+2)=(\frac{9\pi}{8}+2\pi n,1) \] where $n$ is an integer:\newline

\section*{19}
We are given the function $y=3\cos\left(x+\frac{3\pi}{4}\right)-1$.
We know that the horizontal shift is $h=-\frac{3\pi}{4}$ and the vertical shift is $k=-1$. Also, the period is $2\pi$ and the amplitude is $3$. Thus the midline is $y=-1$, and the points on it are
\[ x=\frac{1}{4}\times 2\pi-\frac{3\pi}{4}+2\pi n=-\frac{\pi}{4}+2\pi n\vee x=\frac{3}{4}\times 2\pi-\frac{3\pi}{4}+2\pi n=\frac{3\pi}{4}+2\pi n \] where $n$ is an integer and the maximum 
values are at
\[ (2\pi-\frac{3\pi}{4}+2\pi n,3-1)=(\frac{5\pi}{4}+2\pi n,2) \] where $n$ is an integer and the minimum values are at
\[ (\pi-\frac{3\pi}{4}+2\pi n,-3-1)=(\frac{\pi}{4}+2\pi n,-4) \] where $n$ is an integer:\newline

\section*{20}
We are given the function $y=\sin 2\left(x+2\pi\right)-3$.
We know that the horizontal shift is $h=-2\pi$ and the vertical shift is $k=-3$. Also, the period is $\frac{2\pi}{2}=\pi$ and the amplitude is 1. Thus the midline is $y=-3$, and the points on it are
\[ x=\pi-2\pi+\pi n=-\pi+\pi n\vee x=\frac{\pi}{2}-2\pi+\pi n=-\frac{3\pi}{2}+\pi n \] where $n$ is an integer and the maximum values are at
\[ (\frac{1}{4}\pi-2\pi+\pi n,1-3)=(-\frac{7\pi}{4}+\pi n,-2) \] where $n$ is an integer and the minimum values are at
\[ (\frac{3}{4}\pi-2\pi+\pi n,-1-3)=(-\frac{5\pi}{4}+\pi n,-4) \] where $n$ is an integer:\newline

\section*{21}
C

\end{document}
